% English translation, made from our own transcription (original.tex), not from any existing
% translation. Every math expression, symbol, \tag{}, \origpage{N} marker, and \ednote is
% reproduced unchanged; only the prose is translated — including the short prose set INSIDE math
% via a text insert, which is translated like any other prose (HOUSESTYLE R21): "und"→"and",
% "oder da"→"or since", "ist"→"is", the ordinals "ster"→"th" (so "$m-1\text{ster}$" reads
% "$m-1\text{th}$"), and "w. z. b. w." (was zu beweisen war) → "q. e. d.". pipeline/texcompare.py
% ignores the CONTENT of such an insert while still checking that each insert is present, unmoved,
% and that all surrounding math is identical.
%
% The 12 \ednote notes flagging apparent printer's errors (pp. 9, 11, 13, 17, 18, 20, 24, 25, 31)
% are carried across and translated into English; the errors themselves are in the mathematics
% or typography, so they survive translation unchanged.
%
% TERMINOLOGY. Applied consistently from corpus/riemann-1851-complexen-groesse/glossary.yaml
% and harmonized with riemann-1857-abelsche-functionen:
%   * veränderliche complexe Grösse = "variable complex quantity"
%   * Differentialquotient = "differential quotient"
%   * Grössenoperationen = "operations on quantities"
%   * Grössengebiet = "domain of quantities"
%   * Fläche = "surface", ausgebreitet = "spread out", schlicht = "simple"
%   * zusammenhängend = "connected", einfach zusammenhängend = "simply connected",
%     mehrfach zusammenhängend = "multiply connected", n-fold connected = "n-fold connected"
%   * Ordnung des Zusammenhangs = "order of connectivity"
%   * Queerschnitt = "crosscut", Schnitt = "cut"
%   * Windungspunkt = "winding point", Zweigpunkt = "branch point"
%   * Abbildung = "mapping", in den kleinsten Theilen ähnlich = "similar in the smallest parts"
%   * Lehrsatz = "Theorem", w. z. b. w. / w. z. v. w. = "q. e. d."
\documentclass{article}
\usepackage{readmasters}
\begin{document}

\origpage{1}
\section*{1.}

If one conceives by $z$ a variable quantity which can successively assume all possible real values, then, when to each of its values there corresponds a single value of the indeterminate quantity $w$, $w$ is called a function of $z$; and if, while $z$ continuously passes through all values situated between two fixed values, $w$ likewise varies continuously, then this function is called continuous within this interval.

This definition evidently establishes absolutely no law between the individual values of the function, in that, when this function has been disposed of for a determined interval, the manner of its continuation outside the same remains entirely left to arbitrariness.

The dependence of the quantity $w$ on $z$ may be given by a mathematical law, so that through determined operations on quantities, for each value of $z$ the corresponding $w$ is found. The capacity to be determined for all values of $z$ lying within a given interval by the same law of dependence was formerly ascribed only to a certain genus of functions (\emph{functiones continuae} according to Euler's usage); more recent investigations have shown, however, that there exist analytical expressions by which every continuous function can be represented for a given interval. It is therefore indifferent whether one defines the dependence of the quantity $w$ on the quantity $z$ as arbitrarily given or as conditioned by determined operations on quantities. Both concepts are congruent in consequence of the mentioned theorems.

The matter stands otherwise, however, when the variability of the quantity $z$ is not restricted to real values, but complex ones of the form $x + yi$ (where $i = \sqrt{-1}$) are also admitted.

Let $x + yi$ and $x + yi + dx + dy\,i$ be two infinitely little different values of the quantity $z$, to which correspond the values $u + vi$ and $u + vi + du + dv\,i$ of the quantity $w$. Then, if the dependence of the quantity $w$ on $z$ is an arbitrarily assumed one, the ratio $\frac{du + dv\,i}{dx + dy\,i}$ will, generally speaking, vary with the values of $dx$ and $dy$,

\origpage{2}
since, if one puts $dx + dy\,i = \varepsilon e^{\varphi i}$, it becomes
\[
\begin{aligned}
\frac{du + dv\,i}{dx + dy\,i} &= \frac{1}{2}\left(\frac{du}{dx} + \frac{dv}{dy}\right) + \frac{1}{2}\left(\frac{dv}{dx} - \frac{du}{dy}\right)i + \frac{1}{2}\left[\frac{du}{dx} - \frac{dv}{dy} + \left(\frac{dv}{dx} + \frac{du}{dy}\right)i\right]\frac{dx - dy\,i}{dx + dy\,i} \\
&= \frac{1}{2}\left(\frac{du}{dx} + \frac{dv}{dy}\right) + \frac{1}{2}\left(\frac{dv}{dx} - \frac{du}{dy}\right)i + \frac{1}{2}\left[\frac{du}{dx} - \frac{dv}{dy} + \left(\frac{dv}{dx} + \frac{du}{dy}\right)i\right]e^{-2\varphi i}
\end{aligned}
\]
In whatever way, however, $w$ may be determined as a function of $z$ by compounding the simple operations on quantities, the value of the differential quotient $\frac{dw}{dz}$ will always be independent of the particular value of the differential $dz$${}^{*)}$. Evidently, therefore, not every arbitrary dependence of the complex quantity $w$ on the complex quantity $z$ can be expressed in this way.

The characteristic just highlighted of all functions determinable in any way by operations on quantities we shall lay at the foundation for the following investigation, where such a function is to be considered independently of its expression, by proceeding, without now proving its general validity and sufficiency for the concept of a dependence expressible by operations on quantities, from the following definition:

A variable complex quantity $w$ is called a function of another variable complex quantity $z$ if it varies with it in such a way that the value of the differential quotient $\frac{dw}{dz}$ is independent of the value of the differential $dz$.

\section*{2.}

Both the quantity $z$ and the quantity $w$ are regarded as variable quantities that can assume every complex value. The conception of such a variability, which extends over a connected domain of two dimensions, is substantially facilitated by connecting it with spatial intuitions.

Let each value $x + yi$ of the quantity $z$ be conceived as represented by a point $O$ of the plane $A$, whose rectangular coordinates are $x$, $y$, and each value $u + vi$ of the quantity $w$ by a point $Q$ of the plane $B$, whose rectangular coordinates are $u$, $v$. Every dependence of the quantity $w$ on $z$ will then present itself as a dependence of the position of the point $Q$ on that of the point $O$. If to each value of $z$ there corresponds a determinate value of $w$ varying continuously with $z$, in other words, if $u$ and $v$ are continuous functions of $x$, $y$, then to each point of the plane $A$ there will correspond a point of the plane $B$, to each line, generally speaking, a line, and to each connected piece of surface a connected piece of surface. One will therefore be able to visualize this dependence of the quantity $w$ on $z$ as a mapping of the plane $A$ onto the plane $B$.

\textbf{*)} This assertion is evidently justified in all cases where an expression of $\frac{dw}{dz}$ in terms of $z$ can be found from the expression of $w$ in terms of $z$ by means of the rules of differentiation; its strictly general validity remains for now undecided.

\origpage{3}
\section*{3.}

It shall now be investigated what property this mapping receives when $w$ is a function of the complex quantity $z$, that is, when $\frac{dw}{dz}$ is independent of $dz$.

We denote by $o$ an indeterminate point of the plane $A$ in the vicinity of $O$, its image in the plane $B$ by $q$, and furthermore by $x + yi + dx + dy\,i$ and $u + vi + du + dv\,i$ the values of the quantities $z$ and $w$ at these points. Then $dx$, $dy$ and $du$, $dv$ can be regarded as rectangular coordinates of the points $o$ and $q$ with respect to the points $O$ and $Q$ as origins, and if one puts $dx + dy\,i = \varepsilon e^{\varphi i}$ and $du + dv\,i = \eta e^{\psi i}$, then the quantities $\varepsilon$, $\varphi$, $\eta$, $\psi$ will be polar coordinates of these points for the same origins. If now $o'$ and $o''$ are any two determinate positions of the point $o$ in infinite proximity to $O$, and if one expresses the meanings of the remaining symbols depending on them by corresponding indices, the presupposition gives
\[
\frac{du' + dv'i}{dx' + dy'i} = \frac{du'' + dv''i}{dx'' + dy''i}
\]
and consequently
\[
\frac{du' + dv'i}{du'' + dv''i} = \frac{\eta'}{\eta''} e^{(\psi' - \psi'')i} = \frac{dx' + dy'i}{dx'' + dy''i} = \frac{\varepsilon'}{\varepsilon''} e^{(\varphi' - \varphi'')i},
\]
from which
\[
\frac{\eta'}{\eta''} = \frac{\varepsilon'}{\varepsilon''} \quad \text{and} \quad \psi' - \psi'' = \varphi' - \varphi'',
\]
that is, in the triangles $o'Oo''$ and $q'Qq''$ the angles $o'Oo''$ and $q'Qq''$ are equal and the sides including them are proportional to each other.

There obtains, therefore, similarity between two mutually corresponding infinitely small triangles, and consequently in general between the smallest parts of the plane $A$ and of its image on the plane $B$. An exception to this theorem occurs only in the special cases when the mutually corresponding changes of the quantities $z$ and $w$ do not stand in a finite ratio to each other, which was tacitly presupposed in its derivation${}^{*)}$.

\section*{4.}

If one brings the differential quotient $\frac{du + dv\,i}{dx + dy\,i}$ into the form
\[
\frac{\left(\frac{du}{dx} + \frac{dv}{dx}i\right)dx + \left(\frac{dv}{dy} - \frac{du}{dy}i\right)dy\,i}{dx + dy\,i},
\]

\textbf{*)} On this subject see:
„Allgemeine Auflösung der Aufgabe: die Theile einer gegebenen Fläche so abzubilden, dass die Abbildung dem Abgebildeten in den kleinsten Theilen ähnlich wird, von C.~F.~Gauss. (Als Beantwortung der von der königlichen Societät der Wissenschaften in Copenhagen für 1822 aufgegebenen Preisfrage“, reprinted in: „Astronomische Abhandlungen, herausgegeben von Schumacher. Drittes Heft. Altona. 1825“).

\origpage{4}
it becomes clear that it will have the same value for every two values of $dx$ and $dy$ if and only if
\[
\frac{du}{dx} = \frac{dv}{dy} \quad \text{and} \quad \frac{dv}{dx} = -\frac{du}{dy}
\]
is satisfied. These conditions are therefore sufficient and necessary in order that $w = u + vi$ be a function of $z = x + yi$. For the individual terms of this function there flow from them the following:
\[
\frac{d^2u}{dx^2} + \frac{d^2u}{dy^2} = o, \quad \frac{d^2v}{dx^2} + \frac{d^2v}{dy^2} = o,
\]
which form the foundation for the investigation of the properties that belong to a single term of such a function considered individually. We shall let the proof for the most important of these properties precede a more detailed consideration of the complete function, but first discuss and establish some points belonging to more general domains, in order to level the ground for those investigations.

\section*{5.}

For the following considerations we restrict the variability of the quantities $x$, $y$ to a finite domain, by considering as the locus of the point $O$ no longer the plane $A$ itself, but a surface $T$ spread out over the same. We choose this clothing, in which it will be unexceptionable to speak of surfaces lying upon one another, in order to leave open the possibility that the locus of the point $O$ extends multiple times over the same part of the plane; we presuppose, however, for such a case that the surface parts lying upon one another are not connected along a line, so that a folding over of the surface, or a splitting into parts lying upon one another, does not occur.

The number of surface parts lying upon one another in each part of the plane is then completely determined when the boundary is given in position and sense (that is, its inner and outer side); their course can nevertheless still shape itself differently.

In fact, if we draw an arbitrary line $l$ through the part of the plane covered by the surface, the number of surface parts lying over one another changes only upon crossing the boundary, and indeed upon passing from outside to inside by $+1$, in the opposite case by $-1$, and is thus everywhere determined. Along the shore of this line each adjacent surface part now continues in a completely determined manner, so long as the line does not meet the boundary, since an indeterminacy can in any case take place only at an isolated point and thus either at a point of the line itself or at a finite distance from it; we can therefore, if we restrict our consideration to a part of the line $l$ running in the interior of the surface and to a sufficiently small surface strip on both sides, speak of determinate adjacent surface parts, whose number is the same on each side, and which, by assigning to the line a determinate direction, we denote on the left by $a_1$, $a_2$, \dots\ $a_n$, on the right by $a'_1$, $a'_2$, \dots\ $a'_n$. Each surface part $a$ will then continue into one of the surface parts $a'$; this will indeed in general be the same for the entire course of the line $l$, but can nevertheless change for particular positions of $l$ at one of its points. Let us assume that above such a point $\sigma$ (that is,

\origpage{5}
along the preceding part of $l$) there are connected with the surface parts $a'_{1}$, $a'_{2}$, \dots\ $a'_{n}$ in succession the surface parts $a_{1}$, $a_{2}$, \dots\ $a_{n}$, but below it the surface parts $a_{\alpha_1}$, $a_{\alpha_2}$, \dots\ $a_{\alpha_n}$, where $\alpha_1$, $\alpha_2$, \dots\ $\alpha_n$ differ from $1$, $2$, \dots\ $n$ only in arrangement; then a point entering from $a_1$ into $a'_1$ above $\sigma$, when it steps back onto the left side below $\sigma$, will arrive in the surface part $a_{\alpha_1}$, and when it circles around the point $\sigma$ from left to right, the index of the surface part in which it finds itself will pass in succession through the numbers
\[
1, \, \alpha_1, \, \alpha_{\alpha_1}, \, \dots \, \mu, \, \alpha_\mu, \, \dots
\]
In this series, as long as the term $1$ does not recur, all terms are necessarily different from one another, because an arbitrary middle term $\alpha_\mu$ is necessarily preceded in immediate succession by $\mu$ and successively by all earlier terms down to $1$; but if after a number of terms, which must evidently be less than $n$ and let be $= m$, the term $1$ recurs, then the remaining terms must follow in the same order. The point moving about $\sigma$ then returns after every $m$ circuits into the same surface part and is confined to $m$ of the surface parts lying upon one another, which unite over $\sigma$ into a single point. We call this point a winding point of the $m-1\text{th}$ order of the surface $T$. By applying the same procedure to the remaining $n-m$ surface parts, these, if they do not run separately, will decompose into systems of $m_1$, $m_2$, \dots\ surface parts, in which case there also lie in the point $\sigma$ winding points of the $m_1-1\text{th}$, $m_2-1\text{th}$ order.

When the position and the sense of the boundary of $T$ and the positions of its winding points are given, then $T$ is either completely determined or at least restricted to a finite number of different shapes; the latter, in so far as these determining elements can relate to different ones of the surface parts lying upon one another.

A variable quantity which for each point $O$ of the surface $T$, generally speaking, that is, without excluding an exception in isolated lines and points${}^{*)}$, assumes a single determinate value varying continuously with its position, can evidently be regarded as a function of $x$, $y$, and wherever in the sequel functions of $x$, $y$ are spoken of, we shall establish the concept of the same in this manner.

Before we turn to the consideration of such functions, however, we interpolate some discussions concerning the connectivity of a surface. We restrict ourselves thereby to such surfaces as do not split along a line.

\textbf{*)} This restriction is indeed not required by the concept of a function in itself, but is necessary in order to be able to apply the infinitesimal calculus to it: a function that is discontinuous at all points of a surface, such as e.\ g.\ a function that for a commensurable $x$ and a commensurable $y$ has the value $1$, but otherwise has the value $2$, can be subjected neither to differentiation nor to integration, and thus (directly) not to the infinitesimal calculus at all. The restriction made here arbitrarily for the surface $T$ will justify itself later (Art.~15).

\origpage{6}
\section*{6.}

We consider two surface parts as connected or belonging to one piece if from a point of the one a line can be drawn through the interior of the surface to a point of the other; as separated, if this possibility does not obtain.

The investigation of the connectivity of a surface rests on its decomposition by crosscuts, that is, lines which cut through the interior simply --- no point multiply --- from a boundary point to a boundary point. The latter can also lie in the part added to the boundary, thus in an earlier point of the crosscut.

A connected surface is called, if it falls into pieces by every crosscut, a simply connected one; otherwise a multiply connected one.

\emph{Theorem I.} A simply connected surface $A$ falls by every crosscut $a\,b$ into two simply connected pieces.

Supposing one of these pieces were not broken into pieces by a crosscut $c\,d$, then, according as neither of its endpoints, or the endpoint $c$, or both endpoints fell into $a\,b$, one would evidently obtain, by restoring the connection along the entire line $a\,b$, or along the part $c\,b$, or the part $c\,d$ of the same, a connected surface that would arise from $A$ by a crosscut, contrary to the hypothesis.

\emph{Theorem II.} If a surface $T$ falls by $n_1{}^{*)}$ crosscuts $q_1$ into a system $T_1$ of $m_1$ simply connected surface pieces, and by $n_2$ crosscuts $q_2$ into a system $T_2$ of $m_2$ surface pieces, then $n_2 - m_2$ cannot be $> n_1 - m_1$.

Every line $q_2$, if it does not fall entirely into the system of crosscuts $q_1$, forms at the same time one or more crosscuts $q'_2$ of the surface $T_1$. As endpoints of the crosscuts $q'_2$ are to be regarded:
1) the $2n_2$ endpoints of the crosscuts $q_2$, except when their ends coincide with a part of the line system $q_1$,


2) each middle point of a crosscut $q_2$ at which it enters into a middle point of a line $q_1$, except when it is already in another line $q_1$, that is, when an end of a crosscut $q_1$ coincides with it.
If now $\mu$ denotes how often lines of both systems meet or diverge during their course (where thus an isolated common point is to be counted twice), $\nu_1$ how often an end piece of the $q_1$ coincides with a middle piece of the $q_2$, $\nu_2$ how often an end piece of the $q_2$

\textbf{*)} By a decomposition by several crosscuts is always to be understood a successive one, that is, one where the surface produced by a crosscut is further decomposed by a new crosscut.

\origpage{7}
with a middle piece of the $q_1$, and finally $\nu_3$ how often an end piece of the $q_1$ coincides with an end piece of the $q_2$, then No.\ 1 furnishes $2n_2 - \nu_2 - \nu_3$, No.\ 2 $\mu - \nu_1$ endpoints of the crosscuts $q'_2$; both cases taken together, however, comprise all endpoints and each only once, and the number of these crosscuts is therefore
\[
\frac{2n_2 - \nu_2 - \nu_3 + \mu - \nu_1}{2} = n_2 + s.
\]
By entirely similar inferences there results the number of crosscuts $q'_1$ of the surface $T_2$ formed by the lines $q_1$,
\[
= \frac{2n_1 - \nu_1 - \nu_3 + \mu - \nu_2}{2}, \quad \text{hence} \quad = n_1 + s.
\]
The surface $T_1$ is now evidently transformed by the $n_2 + s$ crosscuts $q'_2$ into the same surface into which $T_2$ is decomposed by the $n_1 + s$ crosscuts $q'_1$. But $T_1$ consists of $m_1$ simply connected pieces and therefore falls according to Theorem I by $n_2 + s$ crosscuts into $m_1 + n_2 + s$ surface pieces; consequently, were $m_2 < m_1 + n_2 - n_1$, the number of surface pieces $T_2$ would have to be increased by $n_1 + s$ crosscuts by more than $n_1 + s$, which is absurd.

According to this theorem, if the number of crosscuts is denoted indeterminately by $n$, and the number of pieces by $m$, then $n - m$ is constant for all decompositions of a surface into simply connected pieces; for if we consider any two determinate decompositions by $n_1$ crosscuts into $m_1$ pieces and by $n_2$ crosscuts into $m_2$ pieces, then, if the former are simply connected, $n_2 - m_2 \le n_1 - m_1$, and if the latter are simply connected, $n_1 - m_1 \le n_2 - m_2$, thus when both obtain, $n_2 - m_2 = n_1 - m_1$.

This number can fitly be designated with the name ``order of connectivity'' of a surface; it is
--- lowered by $1$ by each crosscut --- according to the definition ---,


--- not changed by a line cutting simply through the interior from an interior point to a boundary point or to an earlier point of section, and


--- increased by $1$ by an interior everywhere simple cut terminating in two points,
because the former can be converted by one, but the latter by two crosscuts into one crosscut.

Finally, the order of connectivity of a surface consisting of several pieces is obtained by adding the orders of connectivity of these pieces to one another.

We shall, however, in the sequel mostly restrict ourselves to a surface consisting of one piece, and make use for its connectivity of the simpler designation of a simple, twofold, etc., by understanding by an $n$-fold connected surface one that can be decomposed by $n-1$ crosscuts into a simply connected one.

With regard to the dependence of the connectivity of the boundary on the connectivity of a surface, it is easily seen:

\origpage{8}
1) The boundary of a simply connected surface consists necessarily of a single closed line.
Did the boundary consist of separate pieces, a crosscut $q$ connecting a point of one piece $a$ with a point of another $b$ would separate only connected parts of the surface from one another, since in the interior of the surface along $a$ a line could be led from one side of the crosscut $q$ to the opposite side; and consequently $q$ would not break the surface into pieces, contrary to the hypothesis.

2) By each crosscut the number of boundary pieces is either diminished by $1$ or increased by $1$.
A crosscut $q$ either connects a point of one boundary piece $a$ with a point of another $b$, --- in this case all these lines taken together in the succession $a$, $q$, $b$, $q$ form a single closed piece of the boundary ---

or it connects two points of one piece of the boundary, --- in this case this piece falls by its two endpoints into two pieces, each of which taken together with the crosscut forms a closed boundary piece ---

or finally, it terminates in one of its earlier points and can be regarded as composed of a closed line $o$ and another $l$ connecting a point of $o$ with a point of a boundary piece $a$, --- in which case $o$ on the one hand, and $a$, $l$, $o$, $l$ on the other hand each form a closed boundary piece.

There thus take the place either --- in the first case --- of two pieces, one, or --- in the two latter cases --- of one piece, two boundary pieces, from which our proposition follows.

The number of pieces of which the boundary of an $n$-fold connected piece of surface consists is therefore either $= n$ or smaller by an even number.

From this we further draw the corollary:

If the number of boundary pieces of an $n$-fold connected surface is $= n$, it falls by every everywhere simple closed cut in the interior into two separate pieces.

For the order of connectivity is not changed thereby, while the number of boundary pieces is increased by $2$; the surface, were it a connected one, would therefore have an $n$-fold connectivity and $n+2$ boundary pieces, which is impossible.

\section*{7.}

If $X$ and $Y$ are two functions of $x$, $y$ continuous at all points of the surface $T$ spread out over $A$, the integral extended over all elements $dT$ of this surface is
\[
\int \left(\frac{dX}{dx} + \frac{dY}{dy}\right) dT = -\int (X\cos\xi + Y\cos\eta)\,ds,
\]
when at each point of the boundary the inclination of an inward-drawn normal toward the $x$-axis is denoted by $\xi$, and toward the $y$-axis by $\eta$, and this integration extends over all elements $ds$ of the boundary line.

In order to transform the integral $\displaystyle\int \frac{dX}{dx}\,dT$, we decompose the part of the plane $A$ covered by the surface $T$

\origpage{9}
into elementary strips by a system of lines parallel to the $x$-axis, and indeed in such a way that every winding point of the surface $T$ falls on one of these lines. Under this presupposition, the part of $T$ falling on each of them consists of one or more separately running trapezoidal pieces. The contribution of an indeterminate one of these surface strips, which cuts out of the $y$-axis the element $dy$, to the value of $\displaystyle\int \frac{dX}{dx}\,dT$ will then evidently be $= dy\displaystyle\int \frac{dX}{dx}\,dx$, when this integration is extended through that or those straight lines belonging to the surface $T$ which fall on a normal passing through a point of $dy$. If now the lower endpoints of the same (that is, to which the smallest values of $x$ correspond) are $O_{'}$, $O_{''}$, $O_{'''}$, \dots, the upper ones $O'$, $O''$, $O'''$, \dots, and if we denote by $X_{'}$, $X_{''}$, \dots\ $X'$, $X''$, \dots\ the values of $X$ at these points, by $ds_{'}$, $ds_{''}$, \dots\ $ds'$, $ds''$, \dots\ the corresponding elements cut out of the boundary by the surface strip, and by $\xi_{'}$, $\xi_{''}$, \dots\ $\xi'$, $\xi''$, \dots\ the values of $\xi$ at these elements, then
\[
\begin{aligned}
\int \frac{dX}{dx}\,dx = & - X_{'} - X_{''} - X_{'''} \dots \\
& + X' + X'' + X''' \dots
\end{aligned}
\]
The angles $\xi$ evidently become acute at the lower, obtuse at the upper endpoints, and there will therefore be
\[
\begin{aligned}
dy &= \phantom{-}\cos\xi_{'}\,ds_{'} = \phantom{-}\cos\xi_{''}\,ds_{''} \dots \\
&= -\cos\xi'\,ds' = -\cos\xi''\,ds'' \dots
\end{aligned}
\]
By substitution of these values there results
\[
dy \int \frac{dX}{dx} = -\Sigma X\cos\xi\,ds,
\]
where the summation refers to all boundary elements that have $dy$ as their projection on the $y$-axis.

By integration over all $dy$ coming into consideration, all elements of the surface $T$ and all elements of the boundary are evidently exhausted, and one therefore obtains, taken in this extent,
\[
\int \frac{dX}{dx}\,dT = -\int Y\cos\xi\,ds.
\]\ednote{Printed $Y$ here; mathematically $X$ is required.}
By entirely similar inferences one finds
\[
\int \frac{dY}{dy}\,dT = -\int Y\cos\eta\,ds \quad \text{and consequently}
\]
\[
\int \left(\frac{dX}{dx} + \frac{dY}{dy}\right) dT = -\int (X\cos\xi + Y\cos\eta)\,ds, \quad \text{q.\ e.\ d.}
\]

\section*{8.}

If in the boundary line, reckoned from a fixed initial point in a determinate direction to be fixed later, we denote its length up to an indeter-

\origpage{10}
minate point $O_0$ by $s$, and along the normal erected at this point $O_0$ the distance of an indeterminate point $O$ from the same, considered as positive inwards, by $p$, then the values of $x$ and $y$ at the point $O$ can evidently be regarded as functions of $s$ and $p$, and the partial differential quotients in the points of the boundary line will then be
\[
\frac{dx}{dp} = \cos\xi, \quad \frac{dy}{dp} = \cos\eta, \quad \frac{dx}{ds} = \pm\cos\eta, \quad \frac{dy}{ds} = \mp\cos\xi,
\]
where the upper signs hold when the direction in which the quantity $s$ is regarded as increasing encloses with $p$ the same angle as the $x$-axis with the $y$-axis, and the lower when an opposite one. We shall assume this direction in all parts of the boundary such that $\frac{dx}{ds} = \frac{dy}{dp}$ and consequently $\frac{dy}{ds} = -\frac{dx}{dp}$, which does not essentially impair the generality of our results.

Evidently we can extend these determinations also to lines in the interior of $T$; only here, for the determination of the signs of $dp$ and $ds$, when their mutual dependence is fixed as there, we still have to add a specification fixing either the sign of $dp$ or of $ds$; and indeed, for a closed line we shall specify for which of the surface parts separated by it it is to count as boundary, whereby the sign of $dp$ is determined, but for an unclosed line its initial point, that is, the endpoint where $s$ assumes the smallest value.

The introduction of the values obtained for $\cos\xi$ and $\cos\eta$ into the equation proved in the preceding Art.\ gives, taken in the same extent as there,
\[
\int \left(\frac{dX}{dx} + \frac{dY}{dy}\right) dT = -\int \left(X\frac{dx}{dp} + Y\frac{dy}{dp}\right) ds = \int \left(X\frac{dy}{ds} - Y\frac{dx}{ds}\right) ds.
\]

\section*{9.}

By application of the theorem at the conclusion of the preceding Art.\ to the case where $\frac{dX}{dx} + \frac{dY}{dy} = o$ in all parts of the surface, we obtain the following theorems:

\emph{I.} If $X$ and $Y$ are two functions finite and continuous in all points of $T$ and satisfying the equation $\frac{dX}{dx} + \frac{dY}{dy} = o$, then extended through the whole boundary of $T$
\[
\int \left(X\frac{dx}{dp} + Y\frac{dy}{dp}\right) ds = o.
\]

If one conceives an arbitrary surface $T_1$ spread out over $A$ decomposed in any manner into two pieces $T_2$ and $T_3$, the integral $\displaystyle\int \left(X\frac{dx}{dp} + Y\frac{dy}{dp}\right) ds$ with respect to the boundary of $T_2$ can be regarded as the difference of the integrals with respect to the boundary

\origpage{11}
of $T_1$ and with respect to the boundary of $T_3$, in that, where $T_3$ extends to the boundary of $T_1$, both integrals cancel each other, but all remaining elements correspond to an element of the boundary of $T_2$.

By means of this transformation there results from I.:

\emph{II.} The value of the integral $\displaystyle\int \left(X\frac{dx}{dp} + Y\frac{dy}{dp}\right) ds$, extended through the whole boundary of a surface spread out over $A$, remains constant under arbitrary expansion or contraction of the same, if only no surface parts enter or leave within which the presuppositions of Theorem I are not fulfilled.

If now the functions $X$, $Y$, although satisfying the prescribed differential equation in each part of the surface $T$, are nevertheless affected with a discontinuity in isolated lines or points, one can surround each such line and each such point with an arbitrarily small surface part as envelope, and then obtains by application of Theorem II:

\emph{III.} The integral $\displaystyle\int \left(X\frac{dx}{dp} + Y\frac{dy}{dp}\right) ds$ with respect to the whole boundary of $T$ is equal to the sum of the integrals $\displaystyle\int \left(X\frac{dx}{dp} + Y\frac{dy}{dp}\right) ds$ with respect to the boundaries of all points of discontinuity, and indeed with respect to each individual one of these points the integral retains the same value, within however narrow limits one may enclose it.

This value is necessarily equal to $o$ for a mere point of discontinuity if with the distance $\varrho$ of the point $O$ from the same $\varrho X$ and $\varrho Y$ become infinitely small at the same time; for if with respect to such a point as origin and an arbitrary initial direction one introduces polar coordinates $\varrho$, $\varphi$ and chooses for the boundary a circle described about the same with radius $\varrho$, the integral relating to it is expressed by
\[
\int_o^{2\pi} \left(X\frac{dx}{dp} + Y\frac{dy}{dp}\right) \varrho\,d\varphi
\]
and consequently cannot have a value $x$ different from zero, because, whatever $x$ may be, $\varrho$ can always be assumed so small that apart from sign $\left(X\frac{dx}{dp} - Y\frac{dy}{dp}\right)\varrho$\ednote{Printed $-$ here; mathematically $+$ is intended.} for every value of $\varphi < \frac{x}{2\pi}$, and consequently
\[
\int_0^{2\pi} \left(X\frac{dx}{dp} + Y\frac{dy}{dp}\right) \varrho\,d\varphi < x
\]

\emph{IV.} If in a simply connected surface spread out over $A$, for each surface part the integral $\displaystyle\int \left(X\frac{dx}{dp} + Y\frac{dy}{dp}\right) ds$ or $\displaystyle\int \left(Y\frac{dx}{ds} - X\frac{dy}{ds}\right) ds = o$ extended through its whole boundary, then for any two fixed points $O_0$ and $O$ this integral receives the same value with respect to all lines in it going from $O_0$ to $O$.

Any two lines $s_1$ and $s_2$ connecting the points $O_0$ and $O$ form taken together

\origpage{12}
a closed line $s_3$. This line either itself possesses the property of cutting through no point multiply, or one can decompose it into several everywhere simple closed lines by traversing it from an arbitrary point and separating out, each time one returns to an earlier point, the part traversed in the meantime and considering the following as the immediate continuation of the preceding. Each such line, however, decomposes the surface into a simply and a twofold connected one; it therefore necessarily forms the entire boundary of one of these pieces, and the integral $\displaystyle\int \left(Y\frac{dx}{ds} - X\frac{dy}{ds}\right) ds$ extended through it thus becomes $= o$ according to the hypothesis. The same holds consequently also for the integral extended through the entire line $s_3$, if the quantity $s$ is regarded everywhere as increasing in the same direction; the integrals extended through the lines $s_1$ and $s_2$ must therefore, if this direction remains unchanged, that is, goes in one of them from $O_0$ to $O$ and in the other from $O$ to $O_0$, cancel each other, and thus, if it is reversed in the latter, become equal.

If one now has any arbitrary surface $T$ in which, generally speaking, $\frac{dX}{dx} + \frac{dY}{dy} = o$, one first excludes, if necessary, the points of discontinuity, so that in the remaining surface piece $\displaystyle\int \left(Y\frac{dx}{ds} - X\frac{dy}{ds}\right) ds = o$ for each surface part, and decomposes this by crosscuts into a simply connected surface $T^*$. For every line in the interior of $T^*$ going from a point $O_0$ to another $O$, our integral then has the same value; this value, for which the abbreviation $\displaystyle\int_{O_0}^O \left(Y\frac{dx}{ds} - X\frac{dy}{ds}\right) ds$ may be permitted, is therefore, $O_0$ being thought as fixed, $O$ as movable, a determinate one for each position of $O$ apart from the course of the connecting line, and can consequently be regarded as a function of $x$, $y$. The variation of this function is expressed for a displacement of $O$ along an arbitrary line element $ds$ by $\left(Y\frac{dx}{ds} - X\frac{dy}{ds}\right) ds$, is everywhere continuous in $T^*$, and along a crosscut of $T$ is equal on both sides;

\emph{V.} the integral $Z = \displaystyle\int_{O_0}^O \left(Y\frac{dx}{ds} - X\frac{dy}{ds}\right) ds$ therefore forms, $O_0$ being thought as fixed, a function of $x$, $y$ which varies everywhere continuously in $T^*$, but upon crossing the crosscuts of $T$ changes by a quantity constant along them from one branch point to the other, and of which the partial differential quotient
\[
\frac{dZ}{dx} = Y, \quad \frac{dZ}{dy} = -X \quad \text{is.}
\]

The changes upon crossing the crosscuts depend on a number of mutually independent quantities equal to the number of crosscuts; for if one

\origpage{13}
traverses the system of crosscuts backwards --- the later parts first ---, this variation is everywhere determined when its value is given at the beginning of each crosscut; the latter values, however, are independent of one another.

\section*{10.}

If one puts for the function hitherto denoted by $X$ $u\frac{du'}{dx} - u'\frac{du}{dx}$, and $u\frac{du'}{dy} - u'\frac{du}{dy}$ for $Y$, then
\[
\frac{dX}{dx} + \frac{dY}{dy} = u\left(\frac{d^2u'}{dx^2} + \frac{d^2u'}{dy^2}\right) - u'\left(\frac{d^2u}{dx^2} + \frac{d^2u}{dy^2}\right),
\]
if therefore the functions $u$ and $u'$ satisfy the equations
\[
\frac{d^2u}{dx^2} + \frac{d^2u}{dy^2} = o, \quad \frac{d^2u'}{dx^2} + \frac{d^2u'}{dx^2} = o
\]\ednote{Printed $\frac{d^2u'}{dx^2}$ here; mathematically $\frac{d^2u'}{dy^2}$ is intended.}
then $\frac{dX}{dx} + \frac{dY}{dy} = o$, and to the expression $\displaystyle\int \left(X\frac{dx}{dp} + Y\frac{dy}{dp}\right) ds$, which becomes $= \displaystyle\int \left(u\frac{du'}{dp} - u'\frac{du}{dp}\right) ds$, the theorems of the preceding Art.\ find application.

Let us now make with respect to the function $u$ the presupposition that it, together with its first differential quotients, suffers possible discontinuities in any case not along a line, and that for each point of discontinuity $\varrho\frac{du}{dx}$ and $\varrho\frac{du}{dy}$ become infinitely small at the same time as the distance $\varrho$ of the point $O$ from the same; then the discontinuities of $u$ can remain entirely out of consideration in consequence of the remark to III of the preceding Art.

For then one can in each straight line issuing from a point of discontinuity assume a value $R$ of $\varrho$ such that $\varrho\frac{du}{d\varrho} = \varrho\frac{du}{dx}\frac{dx}{d\varrho} + \varrho\frac{du}{dy}\frac{dy}{d\varrho}$ always remains finite below it, and if $U$ denotes the value of $u$ for $\varrho = R$, and $M$ apart from sign the greatest value of the function $\varrho\frac{du}{d\varrho}$ in that interval, then taken in the same sense there will always be $u - U < M(\log\varrho - \log R)$, and consequently $\varrho(u - U)$, and thus also $\varrho u$, will become infinitely small at the same time as $\varrho$; the same holds according to the hypothesis of $\varrho\frac{du}{dx}$ and $\varrho\frac{du}{dy}$, and consequently, if $u'$ is subject to no discontinuity, also of $\varrho\left(u\frac{du'}{dx} - u'\frac{du}{dx}\right)$ and $\varrho\left(u\frac{du'}{dy} - u'\frac{du}{dy}\right)$; the case discussed in the preceding Art.\ therefore occurs here.

We now further assume that the surface $T$ forming the locus of the point $O$ is spread out everywhere simply over $A$, and we conceive in it an arbitrary fixed point $O_0$ where $u$, $x$, $y$ receive the values $u_0$, $x_0$, $y_0$. The quantity $\frac{1}{2}\log[(x - x_0)^2 + (y - y_0)^2] = \log r$ considered as a function of $x$, $y$ then has the prop-

\origpage{14}
erty that $\frac{d^2\log r}{dx^2} + \frac{d^2\log r}{dy^2} = o$, and is affected with a discontinuity only for $x = x_0$, $y = y_0$, thus in our case only for a single point of the surface $T$.

There is therefore according to Art.~9, III, if we put $\log r$ for $u'$, $\displaystyle\int \left(u\frac{d\log r}{dp} - \log r\frac{du}{dp}\right) ds$ with respect to the whole boundary of $T$ equal to this integral with respect to an arbitrary boundary of the point $O_0$, and thus, if for this we choose the circumference of a circle where $r$ has a constant value, and denote from one of its points in an arbitrary fixed direction the arc up to $O$ in parts of the radius by $\varphi$, equal to
\[
-\int_0^{2\pi} u\frac{d\log r}{dr} r\,d\varphi - \log r\int \frac{du}{dp}\,ds, \quad \text{or since} \quad \int \frac{du}{dp}\,ds = o \quad \text{is,} \quad = -\int_0^{2\pi} u\,d\varphi,
\]
which value, if $u$ is continuous at the point $O_0$, passes over for an infinitely small $r$ into $-u_0 2\pi$.

Under the presuppositions made with respect to $u$ and $T$ we therefore have for an arbitrary point $O_0$ in the interior of the surface in which $u$ is continuous,
\[
\begin{aligned}
u_0 &= \frac{1}{2\pi}\int \left(\log r\frac{du}{dp} - u\frac{d\log r}{dp}\right) ds \quad \text{with respect to the whole boundary of the same, and} \\
&= \frac{1}{2\pi}\int_0^{2\pi} u\,d\varphi \quad \text{with respect to a circle described about } O_0 \text{.}
\end{aligned}
\]

From the first of these expressions we draw the following

\emph{Theorem.} If a function $u$ within a surface $T$ covering the plane $A$ everywhere simply satisfies, generally speaking, the differential equation $\frac{d^2u}{dx^2} + \frac{d^2u}{dy^2} = o$, and indeed in such a way that
1) the points in which this differential equation is not fulfilled form no surface part,


2) the points in which $u$, $\frac{du}{dx}$, $\frac{du}{dy}$ become discontinuous continuously fill no line,


3) for each point of discontinuity the quantities $\varrho\frac{du}{dx}$, $\varrho\frac{du}{dy}$ become infinitely small at the same time as the distance $\varrho$ of the point $O$ from the same, and


4) for $u$ a discontinuity removable by altering its value in isolated points is excluded,
then it is necessarily finite and continuous, together with all its differential quotients, for all points in the interior of this surface.

In fact, if we consider the point $O_0$ as movable, there vary in the expression $\displaystyle\int \left(\log r\frac{du}{dp} - u\frac{d\log r}{dp}\right) ds$ only the values $\log r$, $\frac{d\log r}{dx}$, $\frac{d\log r}{dy}$. These quantities, however, for each element of the boundary, as long as $O_0$ remains in the interior of $T$, are together with all

\origpage{15}
their differential quotients finite and continuous functions of $x_0$, $y_0$, since the differential quotients are expressed by fractional rational functions of these quantities containing only powers of $r$ in the denominator. The same holds therefore also for the value of our integral, and consequently for the function $u_0$. For this could under the former presuppositions have a value different from it only in isolated points, by becoming discontinuous, which possibility is removed by presupposition 4 of our theorem.

\section*{11.}

Under the same presuppositions with respect to $u$ and $T$ as at the conclusion of the preceding Art., we have the following theorems:

\emph{I.} If along a line $u = o$ and $\frac{du}{dp} = o$, then $u$ is everywhere $= o$.

We prove first that a line $\lambda$ where $u = o$ and $\frac{du}{dp} = o$ cannot form the boundary of a surface part $a$ where $u$ is positive.

Supposing this took place, one would separate from $a$ a piece which is bounded on the one hand by $\lambda$, on the other hand by a circular line, and does not contain the center of this circle, which construction is always possible. One then has, if one denotes the polar coordinates of $O$ with respect to $O_0$ by $r$, $\varphi$, extended through the whole boundary of this piece $\displaystyle\int \log r\frac{du}{dp}\,ds - \displaystyle\int u\frac{d\log r}{dp}\,ds = o$, thus in consequence of the assumption also for the entire circular arc belonging to it
\[
\int u\,d\varphi + \log r \int \frac{du}{dp}\,ds = o, \quad \text{or since} \quad \int \frac{du}{dp}\,ds = o \quad \text{is,} \quad \int u\,d\varphi = o,
\]
which is incompatible with the presupposition that $u$ is positive in the interior of $a$.

In a similar manner it is proved that the equations $u = o$ and $\frac{du}{dp} = o$ cannot take place in a boundary part of a surface piece $b$ where $u$ is negative.

If now in the surface $T$ in a line $u = o$ and $\frac{du}{dp} = o$, and in any part of it $u$ were different from zero, such a surface part would evidently have to be bounded either by this line itself or by a surface part where $u = o$, thus in any case by a line where $u$ and $\frac{du}{dp} = o$, which leads necessarily to one of the assumptions refuted above.

\emph{II.} If the value of $u$ and $\frac{du}{dp}$ is given along a line, then $u$ is thereby determined in all parts of $T$.

If $u_1$ and $u_2$ are any two determinate functions satisfying the

\origpage{16}
conditions imposed on the function $u$, then this holds also, as appears immediately by substitution into these conditions, for their difference $u_1 - u_2$. If now $u_1$ and $u_2$ agreed along a line together with their first differential quotients with respect to $p$, but not in another surface part, there would be along this line $u_1 - u_2 = o$ and $\frac{d(u_1 - u_2)}{dp} = o$, without being everywhere $= o$, contrary to Theorem I.

\emph{III.} The points in the interior of $T$ where $u$ has a constant value form, if $u$ is not everywhere constant, necessarily lines which separate surface parts where $u$ is greater from surface parts where $u$ is smaller.

This theorem is composed of the following:

$u$ cannot have a minimum or a maximum at a point in the interior of $T$;

$u$ cannot be constant \emph{only} in a part of the surface;

the lines in which $u = a$ cannot bound on both sides surface parts where $u - a$ has the same sign;

theorems whose opposite, as is easily seen, would always have to bring about a violation of the equation proved in the preceding Art.
\[
u_0 = \frac{1}{2\pi}\int_0^{2\pi} u\,d\varphi \quad \text{or} \quad \int_0^{2\pi} (u - u_0)\,d\varphi = o
\]
and is consequently impossible.

\section*{12.}

We now turn back to the consideration of a variable complex quantity $w = u + vi$ which, generally speaking (that is, without excluding an exception in isolated lines and points), has for each point $O$ of the surface $T$ a single determinate value varying continuously with its position and in accordance with the equations $\frac{du}{dx} = \frac{dv}{dy}$, $\frac{du}{dy} = -\frac{dv}{dx}$, and we denote this property of $w$, according to what was established earlier, by calling $w$ a function of $z = x + yi$. For the simplification of what follows we stipulate in advance that in a function of $z$ a discontinuity removable by altering its value at an isolated point shall not occur.

The surface $T$ is for the present assigned a simple connectivity and an everywhere simple spread over the plane $A$.

\emph{Theorem.} If a function $w$ of $z$ suffers an interruption of continuity in any case not along a line, and furthermore for each arbitrary point $O'$ of the surface where $z = z'$, $w(z - z')$ becomes infinitely small with the infinite approach of the point $O$, then it is necessarily finite and continuous, together with all its differential quotients, at all points in the interior of the surface.

The presuppositions made concerning the variations of the quantity $w$ resolve, if $z - z' = \varrho e^{\varphi i}$ is put, for $u$ and $v$ into the following: 1) $\frac{du}{dx} - \frac{dv}{dy} = o$ and 2) $\frac{du}{dy} + \frac{dv}{dx} = o$ for every part of the surface $T$; 3) the functions $u$ and $v$ are not discontinuous along a

\origpage{17}
line; 4) for each point $O'$ there become infinitely small with the distance $\varrho$ of the point $O$ from the same $\varrho u$ and $\varrho v$; 5) for the functions $u$ and $v$ discontinuities that could be removed by altering their value at isolated points are excluded.

In consequence of presuppositions 2, 3, 4, for every part of the surface $T$ the integral $\displaystyle\int \Bigl(u\,\frac{dx}{ds} - v\,\frac{dy}{ds}\Bigr)\,ds$ extended over its whole boundary is according to Art.~9, III $= o$, and the integral $\displaystyle\int_{O_0}^O \Bigl(\frac{dx}{ds} - v\,\frac{dy}{ds}\Bigr)\,ds$\ednote{In the print $u$ is missing before $dx/ds$; $u\,dx/ds - v\,dy/ds$ is meant.} therefore receives (according to Article 9, IV) the same value extended through every line going from $O_0$ to $O$, and forms, $O_0$ being thought as fixed, a function $U$ of $x$, $y$ necessarily continuous up to isolated points, of which (and indeed according to 5, at each point) the differential quotient is $\frac{dU}{dx} = u$ and $\frac{dU}{dy} = -v$. By substitution of these values for $u$ and $v$, however, the presuppositions 1, 3, 4 pass over into the conditions of the theorem at the conclusion of Art.~10. The function $U$ is therefore finite and continuous, together with all its differential quotients, at all points of $T$, and the same holds consequently also for the complex function $w = \frac{dU}{dx} - \frac{dU}{dy}\,i$ and its differential quotients taken with respect to $z$.

\section*{13.}

It shall now be investigated what occurs if, while maintaining the other presuppositions of Art.~12, we assume that for a determinate point $O'$ in the interior of the surface, $(z - z')\,w = \varrho e^{\varphi i}w$ no longer becomes infinitely small upon the infinite approach of the point $O$. In this case, then, $w$ becomes infinitely large upon the infinite approach of the point $O$ to $O'$, and we assume that, if the quantity $w$ does not remain of the same order as $\frac{1}{\varrho}$, that is, the quotient of the two approaches a finite limit, at least the orders of both quantities stand in a finite ratio to each other, so that a power of $\varrho$ can be specified whose product with $w$ for an infinitely small $\varrho$ either becomes infinitely small or remains finite. If $\mu$ is the exponent of such a power and $n$ the next greater integer, then the quantity $(z - z')^n w = \varrho^n e^{n\varphi i} w$ becomes infinitely small with $\varrho$, and $(z - z')^{n-1} w$ is therefore a function of $z$ (since $\frac{d(z - z')^{n-1} w}{dz}$ is independent of $dz$) which in this part of the surface satisfies the presuppositions of Art.~12 and is consequently finite and continuous at the point $O'$. If we denote its value at the point $O'$ by $a_{n-1}$, then $(z - z')^{n-1} w - a_{n-1}$ is a function which is continuous and $= o$ at this point and consequently becomes infinitely small with $\varrho$, from which one concludes according to Article 12 that

\origpage{18}
$(z - z')^{n-1}w - \frac{a_{n-1}}{z - z'}$\ednote{Thus in the print; $(z - z')^{n-2}w - a_{n-1}/(z - z')$ is meant.} is a function continuous at the point $O'$. By continuing this procedure, $w$ is evidently converted, by means of subtraction of an expression of the form
\[
\frac{a_1}{z - z'} + \frac{a_2}{(z - z')^2} \dots \frac{a_{n-1}}{(z - z')^{n-1}}
\]
into a function that remains finite and continuous at the point $O'$.

If therefore under the presuppositions of Art.~12 the change occurs that upon the infinite approach of $O$ to a point $O'$ in the interior of the surface $T$ the function $w$ becomes infinitely large, the order of this infinitely large (a quantity growing in the inverse ratio of the distance being regarded as an infinitely large of the first order), if it is finite, is necessarily an integer; and if this number is $= m$, the function $w$ can be converted into one continuous at this point $O'$ by addition of a function containing $2m$ arbitrary constants.

\emph{Note.} We consider a function as containing one arbitrary constant if the possible ways of determining it comprise a continuous domain of one dimension.

\section*{14.}

The restrictions made in Art.~12 and 13 with respect to the surface $T$ are not essential for the validity of the results obtained. Evidently one can surround each point in the interior of an arbitrary surface with a piece of the same possessing the properties there presupposed, with the sole exception of the case where this point is a winding point of the surface.

In order to investigate this case, we conceive the surface $T$ or an arbitrary piece of the same containing a winding point of the $(n-1)$th order $O'$ where let $z = z' = x' = y'i$\ednote{Thus in the print; $z = z' = x' + y'i$ is meant.}, as mapped by means of the function $\zeta = (z - z')^{\frac{1}{n}}$ onto another plane $A$, that is, we conceive the value of the function $\zeta = \xi + \eta i$ at the point $O$ represented by a point $\varTheta$ whose rectangular coordinates are $\xi$, $\eta$ in this plane, and consider $\varTheta$ as the image of the point $O$. In this way one obtains as the mapping of this part of the surface $T$ a connected surface spread out over $\lambda$\ednote{Thus in the print; $A$ is meant.}, which has no winding point at the point $\varTheta'$, the image of the point $O'$, as shall be shown immediately.

To fix ideas, let one conceive described about the point $O'$ in the plane $A$ with radius $R$ a circle, and drawn parallel to the $x$-axis a diameter, where thus $z - z'$ will assume real values. The piece of the surface $T$ surrounding the winding point cut out by this circle will then fall on both sides of the diameter into $n$ separately running semicircular surface pieces, if $R$ is chosen sufficiently small. We denote on that side of the diameter where $y - y'$ is positive these surface pieces by $a_1$, $a_2$ \dots\ $a_n$, on the opposite side by $a'_1$, $a'_2$ \dots\ $a'_n$, and assume that for negative values of $z - z'$, $a_1$, $a_2$ \dots\ $a_n$ are connected in

\origpage{19}
succession with $a'_1$, $a'_2$ \dots\ $a'_n$, for positive on the contrary with $a'_n$, $a'_1$ \dots\ $a'_{n-1}$, so that a point circling the point $O'$ (in the required sense) passes in succession through the surfaces $a_1$, $a'_1$, $a_2$, $a'_2$ \dots\ $a_n$, $a'_n$ and returns through $a'_n$ again into $a_1$, which assumption is evidently permissible. If we now introduce polar coordinates for both planes by putting $z - z' = \varrho e^{\varphi i}$, $\zeta = \sigma e^{\psi i}$, and choose for the mapping of the surface piece $a_1$ that value of $(z - z')^{\frac{1}{n}} = \varrho^{\frac{1}{n}} e^{\frac{1}{n}\varphi i}$ which the latter expression receives under the assumption $o \leqq \varphi \leqq \pi$, then for all points of $a_1$ $\sigma \leqq R^{\frac{1}{n}}$ and $o \leqq \psi \leqq \frac{\pi}{n}$; the images of the same in the plane $A$ thus fall all into a sector, extending from $\psi = o$ to $\psi = \frac{\pi}{n}$, of a circle described about $\varTheta'$ with radius $R^{\frac{1}{n}}$, and indeed to each point of $a_1$ there corresponds one point of this sector advancing continuously with the same and conversely, from which it follows that the mapping of the surface $a_1$ is a connected surface simply spread out over this sector. In a similar manner one obtains for the surface $a'_1$ as mapping a sector extending from $\psi = \frac{\pi}{n}$ to $\psi = \frac{2\pi}{n}$, for $a_2$ one extending from $\psi = \frac{2\pi}{n}$ to $\psi = \frac{3\pi}{n}$, and finally for $a'_n$ one extending from $\psi = \frac{2n - 1}{n}\pi$ to $\psi = 2\pi$, if one chooses $\varphi$ for each point of these surfaces in succession between $\pi$ and $2\pi$, $2\pi$ and $3\pi$ \dots\ $(2n-1)\pi$ and $2n\pi$, which is always and only in one way possible. These sectors, however, join to one another in the same sequence as the surfaces $a$ and $a'$, and indeed in such a way that to the points meeting here there also correspond points meeting there; they can therefore be joined together into a connected mapping of a piece of the surface $T$ enclosing the point $O'$, and this mapping is evidently a surface simply spread out over the plane $A$.

A variable quantity that has a determinate value for each point $O$ has this also for each point $\varTheta$ and conversely, since to each $O$ corresponds only one $\varTheta$ and to each $\varTheta$ only one $O$; if it is furthermore a function of $z$, it is so also of $\zeta$, in that, if $\frac{dw}{dz}$ is independent of $dz$, $\frac{dw}{d\zeta}$ is also independent of $d\zeta$, and conversely. It results from this that to all functions $w$ of $z$ the theorems of Art.~12 and 13 can be applied also at the winding point $O'$, if one considers them as functions of $(z - z')^{\frac{1}{n}}$. This furnishes the following theorem:

If a function $w$ of $z$ becomes infinite upon the infinite approach of $O$ to a winding point of the $(n-1)$th order $O'$, this infinitely large is necessarily of the same

\origpage{20}
order as a power of the distance whose exponent is a multiple of $\frac{1}{n}$, and can, if this exponent is $= -\frac{m}{n}$, be converted into one continuous at the point $O'$ by addition of an expression of the form
\[
\frac{a_1}{(z - z')^{\frac{1}{n}}} + \frac{a_2}{(z - z')^{\frac{2}{n}}} \dots \frac{a_m}{(z - z')^{\frac{m}{n}}},
\]
where $a_1$, $a_2$ \dots\ $a_m$ are arbitrary complex quantities.

This theorem contains as a corollary that the function $w$ is continuous at the point $O'$ if $(z - z')^{\frac{1}{n}} w$ becomes infinitely small upon the infinite approach of the point $O$ to $O'$.

\section*{15.}

Let us now conceive a function of $z$, which has a determinate value for each point $O$ of the surface $T$ arbitrarily spread out over $A$ and is not everywhere constant, represented geometrically, so that its value $w = u + vi$ at the point $O$ is represented by a point $Q$ of the plane $B$ whose rectangular coordinates are $u$, $v$; then there results the following:

\emph{I.} The aggregate of points $Q$ can be regarded as forming a surface $S$ in which to each point there corresponds a single determinate point $O$ advancing continuously with it in $T$.

To prove this, there is evidently required only the verification that the position of the point $Q$ always (and indeed, generally speaking, continuously) varies with that of the point $O$. This is contained in the theorem:

A function $w = u + vi$ of $z$ cannot be constant along a line if it is not everywhere constant. Proof: If $w$ had along a line a constant value $a + bi$, then $u - a$ and $\frac{d(u - a)}{dp}$, which is $= \frac{dv}{ds}$, would be for this line, and $\frac{d^2(u - a)}{dx^2} + \frac{d^2(u - a)}{dy^2}$ everywhere, $= o$; there would therefore have to be according to Theorem 11, I $u - a$, and consequently since $\frac{du}{dx} = \frac{dv}{dy}$, $\frac{du}{dy} = -\frac{dv}{dx}$, also $v - b$, everywhere $= o$, contrary to the hypothesis.

\emph{II.} In consequence of the presupposition made in I, a connection between the parts of $S$ cannot take place without a connection of the corresponding parts of $T$; conversely, wherever in $T$ connection takes place and $w$ is continuous, a corresponding connection can be assigned to the surface $S$.

This being presupposed, the boundary of $S$ corresponds on the one hand to the boundary of $T$, on the other hand to the points of discontinuity; its interior parts, however, isolated points excepted, are everywhere simply spread out over $B$, that is, nowhere does a splitting into parts lying upon one another and nowhere a folding over take place.

The former, since $T$ everywhere possesses a corresponding connection, could evidently occur only if a splitting occurred in $T$ --- contrary to the assumption ---; the latter shall be proved immediately.

\origpage{21}
We prove first of all that a point $Q'$ where $\frac{dw}{dz}$ is finite cannot lie in a fold of the surface $S$.

In fact, if we surround the point $O'$ corresponding to $Q'$ with a piece of the surface $T$ of arbitrary shape and indeterminate dimensions, its dimensions must (according to Art.~3) always be able to be assumed so small that the shape of the corresponding part of $S$ differs arbitrarily little, and consequently so small that its boundary cuts out of the plane $B$ a piece enclosing $Q'$. But this is impossible if $Q'$ lies in a fold of the surface $S$.

Now $\frac{dw}{dz}$, as a function of $z$, can according to I become $= o$ only in isolated points, and, since $w$ is continuous at the points of $T$ coming into consideration, infinite only at the winding points of this surface; consequently etc.\ q.\ e.\ d.

3.\ednote{Thus in the print; in continuation of I and II, III would be expected.} The surface $S$ is consequently a surface for which the presuppositions made in Art.~5 for $T$ hold; and in this surface the indeterminate quantity $z$ has for each point $Q$ a single determinate value which varies continuously with the position of $Q$ and in such a way that $\frac{dz}{dw}$ is independent of the direction of the change of position. In the sense established earlier, $z$ therefore forms a continuous function of the variable complex quantity $w$ for the domain of quantities represented by $S$.

From this it follows further:

If $O'$ and $Q'$ are two corresponding interior points of the surfaces $T$ and $S$, and in the same $z = z'$, $w = w'$, then, if neither of them is a winding point, $\frac{w - w'}{z - z'}$ approaches a finite limit upon the infinite approach of $O$ to $O'$, and the mapping is there one similar in the smallest parts; but if $Q'$ is a winding point of the $n-1$th, $O'$ a winding point of the $m-1$th order, then $\frac{(w - w')^{\frac{1}{n}}}{(z - z')^{\frac{1}{m}}}$ approaches a finite limit upon the infinite approach of $O$ to $O'$, and for the adjoining surface parts there obtains a mode of mapping that easily results from Art.~14.

\section*{16.}

\emph{Theorem.} If $\alpha$ and $\beta$ are two arbitrary functions of $x$, $y$ for which the integral
\[
\int \biggl[\Bigl(\frac{d\alpha}{dx} - \frac{d\beta}{dy}\Bigr)^2 + \Bigl(\frac{d\alpha}{dy} + \frac{d\beta}{dx}\Bigr)^2\biggr]\,dT
\]
extended through all parts of the surface $T$ arbitrarily spread out over $A$ has a finite value, then upon varying $\alpha$ by functions continuous or at least discontinuous only at isolated points which are $= o$ at the boundary, the integral

\origpage{22}
always receives a minimum value for one of these functions, and, if one excludes discontinuities removable by alteration at isolated points, for only one.

We denote by $\lambda$ an indeterminate function, continuous or at least discontinuous only at isolated points, which is $= o$ at the boundary and for which the integral
\[
L = \int \biggl[\Bigl(\frac{d\lambda}{dx}\Bigr)^2 + \Bigl(\frac{d\lambda}{dy}\Bigr)^2\biggr] dT
\]
extended over the whole surface receives a finite value, by $\omega$ an indeterminate one of the functions $\alpha + \lambda$, and finally the integral $\displaystyle\int \biggl[\Bigl(\frac{d\omega}{dx} - \frac{d\beta}{dy}\Bigr)^2 + \Bigl(\frac{d\omega}{dy} + \frac{d\beta}{dx}\Bigr)^2\biggr] dT$ extended over the whole surface by $\varOmega$. The totality of the function $\lambda$ forms a connected domain closed in itself, in that each of these functions can pass continuously into every other, but cannot infinitely approach one discontinuous along a line without $L$ becoming infinite (Art.~16.); for each $\lambda$, setting $\omega = \alpha + \lambda$, $\varOmega$ now receives a finite value which becomes infinite together with $L$, varies continuously with the form of $\lambda$, but can never sink below zero; consequently $\varOmega$ has a minimum for at least one form of the function $\omega$.

In order to prove the second part of our theorem, let $u$ be one of the functions $\omega$ that imparts to $\varOmega$ a minimum value, $h$ an indeterminate quantity constant throughout the entire surface, so that $u + h\lambda$ satisfies the conditions prescribed for the function $\omega$. The value of $\varOmega$ for $\omega = u + h\lambda$, which
\[
\begin{aligned}
&= \int \biggl[\Bigl(\frac{du}{dx} - \frac{d\beta}{dy}\Bigr)^2 + \Bigl(\frac{du}{dy} + \frac{d\beta}{dx}\Bigr)^2\biggr] dT \\
&\quad + 2h \int \biggl[\Bigl(\frac{du}{dx} - \frac{d\beta}{dy}\Bigr)\frac{d\lambda}{dx} + \Bigl(\frac{du}{dy} + \frac{d\beta}{dx}\Bigr)\frac{d\lambda}{dy}\biggr] dT \\
&\quad + h^2 \int \biggl[\Bigl(\frac{d\lambda}{dx}\Bigr)^2 + \Bigl(\frac{d\lambda}{dy}\Bigr)^2\biggr] dT = M + 2Nh + Lh^2 \text{ becomes},
\end{aligned}
\]
must then for each $\lambda$ (according to the concept of the minimum) become greater than $M$, as soon as $h$ is only taken sufficiently small. This requires, however, that for each $\lambda$ $N = o$; for otherwise $2Nh + Lh^2 = Lh^2\bigl(1 + \frac{2N}{Lh}\bigr)$ would become negative if $h$ were assumed opposite to $N$ and apart from sign $< \frac{2N}{L}$. The value of $\varOmega$ for $\omega = u + \lambda$, in which form all possible values of $\omega$ are evidently contained, therefore becomes $= M + L$, and consequently, since $L$ is essentially positive, $\varOmega$ can receive for no form of the function $\omega$ a smaller value than for $\omega = u$.

If now there obtains for another $u'$ of the functions $\omega$ a minimum value $M'$ of $\varOmega$, the same must evidently hold of this, so one has $M' \leqq M$ and $M \leqq M'$, consequently $M = M'$. If one brings $u'$, however, into the form $u + \lambda'$, one obtains for $M'$ the expression $M + L'$, if $L'$ denotes the value of $L$ for $\lambda = \lambda'$, and the equation $M = M'$ gives $L' = o$. This is possible only if in all surface parts $\frac{d\lambda'}{dx} = o$, $\frac{d\lambda'}{dy} = o$, and this function therefore necessarily has, as far as $\lambda'$ is continuous,

\origpage{23}
a constant value, and consequently, since it is $= o$ at the boundary and not discontinuous along a line, at most at isolated points a value different from zero. Two of the functions $\omega$ that impart to $\varOmega$ a minimum value can therefore differ from each other only at isolated points, and if in the function $u$ all discontinuities removable by alteration at isolated points are eliminated, it is completely determined.

\section*{17.}

The proof shall now be supplied that $\lambda$ cannot, without detriment to the finiteness of $L$, infinitely approach a function $\gamma$ discontinuous along a line, that is, if the function $\lambda$ is subjected to the condition of agreeing with $\gamma$ outside a surface part $T'$ enclosing the line of discontinuity, $T'$ can always be assumed so small that $L$ must become greater than an arbitrarily given quantity $C$.

We denote, $s$ and $p$ with respect to the line of discontinuity being taken in the customary meaning, for an indeterminate $s$ the curvature, one convex on the side of positive $p$ being regarded as positive, by $\varkappa$, the value of $p$ at the boundary of $T'$ on the positive side by $p_1$, on the negative side by $p_2$, and the corresponding values of $\gamma$ by $\gamma_1$ and $\gamma_2$. If we now consider any continuously curved part of this line, the part of $T'$ contained between the normals at the endpoints, if it does not extend to the centers of curvature, delivers to $L$ the contribution
\[
\int ds \int_{p_2}^{p_1} dp\,(1 - \varkappa p) \biggl[\Bigl(\frac{d\lambda}{dp}\Bigr)^2 + \Bigl(\frac{d\lambda}{ds}\Bigr)^2 \frac{1}{(1 - \varkappa p)^2}\biggr];
\]
the smallest value of the expression
\[
\int_{p_2}^{p_1} \Bigl(\frac{d\lambda}{dp}\Bigr)^2 (1 - \varkappa p)\,dp
\]
for the fixed boundary values $\gamma_1$ and $\gamma_2$ of $\lambda$ is found, however, according to known rules to be
\[
= \frac{(\gamma_1 - \gamma_2)^2 \varkappa}{\log(1 - \varkappa p_2) - \log(1 - \varkappa p_1)}
\]
and consequently that contribution will necessarily be, however $\lambda$ may be assumed within $T'$,
\[
> \int \frac{(\gamma_1 - \gamma_2)^2 \varkappa\,ds}{\log(1 - \varkappa p_2) - \log(1 - \varkappa p_1)}.
\]
The function $\gamma$ would be continuous for $p = o$ if the greatest value that $(\gamma_1 - \gamma_2)^2$ can receive for $\pi_1 > p_1 > o$ and $\pi_2 < p_2 < o$ became infinitely small with $\pi_1 - \pi_2$; we can consequently for each value of $s$ assume a finite quantity $m$ such that, however small $\pi_1 - \pi_2$ may be assumed, there are always contained within the limits expressed by $\pi_1 > p_1 \geqq o$ and $\pi_2 < p_2 \leqq o$ (where the equalities exclude each other) values of $p_1$ and $p_2$ for which $(\gamma_1 - \gamma_2)^2 > m$. Let us further assume arbitrarily, under the former restrictions, a shape of $T'$ by assigning to $p_1$ and $p_2$ determinate values $P_1$ and $P_2$, and denote the value of the integral

\origpage{24}
extended through the part of the line of discontinuity taken into consideration
\[
\int \frac{m\varkappa\,ds}{\log(1 - \varkappa P_2) - \log(1 - \varkappa P_1)} \quad \text{by } a\text{, then we can evidently make}
\]
\[
\int \frac{(\gamma_1 - \gamma_2)^2 \varkappa\,ds}{\log(1 - \varkappa p_2) - \log(1 - \varkappa p_1)} > C
\]
by assuming $p_1$ and $p_2$ for each value of $s$ such that the inequalities
\[
p_1 < \frac{1 - (1 - \varkappa P_1)^{\frac{a}{C}}}{\varkappa}, \quad p_2 > \frac{1 - (1 - \varkappa P_2)^{\frac{a}{C}}}{\varkappa} \quad \text{and} \quad (\gamma_1 - \gamma_2)^2 > m
\]
are satisfied. This has the consequence, however, that, however $L$ may be assumed within $T'$, the part of $L$ originating from the piece of $T'$ taken into consideration, and consequently all the more $L$ itself, becomes $> C$, q.\ e.\ d.\ednote{Thus in the print; presumably w.\ z.\ b.\ w. [q.\ e.\ d.] is meant.}

\section*{18.}

According to Art.~16 we have for the function $u$ there determined and for any of the functions $\lambda$
\[
N = \int \biggl[\Bigl(\frac{du}{dx} - \frac{d\beta}{dy}\Bigr)\frac{d\lambda}{dx} + \Bigl(\frac{du}{dy} + \frac{d\beta}{dx}\Bigr)\frac{d\lambda}{dy}\biggr] dT
\]
extended through the entire surface $T$ equal to $= o$. From this equation further inferences shall now be drawn.

If one separates from the surface $T$ a piece $T'$ enclosing the points of discontinuity of $u$, $\beta$, $\lambda$, the part of $N$ originating from the remaining piece $T''$ is found with the aid of Art.~7, if one puts $\Bigl(\frac{du}{dx} - \frac{d\beta}{dy}\Bigr)\lambda$ for $X$ and $\Bigl(\frac{du}{dy} + \frac{d\beta}{dx}\Bigr)\lambda$ for $Y$,
\[
= -\int \lambda \Bigl(\frac{d^2u}{dx^2} + \frac{d^2u}{dy^2}\Bigr) dT - \int \Bigl(\frac{du}{dp} - \frac{d\beta}{ds}\Bigr)\lambda\,ds.
\]
In consequence of the boundary condition imposed on the function $\lambda$, the part of $\displaystyle\int \Bigl(\frac{du}{dp} - \frac{d\beta}{ds}\Bigr)\lambda\,ds$ relating to the boundary piece of $T''$ held in common with $T$ becomes equal to $o$, so that $N$ can be regarded as composed of the integral
\[
-\int \lambda \Bigl(\frac{d^2u}{dx^2} + \frac{d^2u}{dy^2}\Bigr) dT \quad \text{with respect to } T'' \quad \text{and}
\]
\[
\int \biggl[\Bigl(\frac{du}{dx} - \frac{d\beta}{dy}\Bigr)\frac{d\lambda}{dx} + \Bigl(\frac{du}{dy} + \frac{d\beta}{dx}\Bigr)\frac{d\lambda}{dy}\biggr] dT + \int \Bigl(\frac{du}{dp} - \frac{d\beta}{ds}\Bigr)\lambda\,ds \quad \text{with respect to } T'.
\]

Evidently now, if $\frac{d^2u}{dx^2} + \frac{d^2u}{dy^2}$ were different from $o$ in any part of the surface $T$, $N$ would likewise receive a value different from $o$ as soon as one chose $\lambda$, as is permitted, equal to $o$ within $T'$ and within $T''$ such that $\lambda \Bigl(\frac{d^2u}{dx^2} + \frac{d^2u}{dy^2}\Bigr)$ had everywhere the same sign. But if $\frac{d^2u}{dx^2} + \frac{d^2u}{dy^2}$ is in all parts of $T = o$, the component

\origpage{25}
of $N$ originating from $T''$ vanishes for each $\lambda$, and the condition $N = o$ then yields that the components relating to the points of discontinuity become $= o$.

For the functions $\frac{du}{dx} - \frac{d\beta}{dy}$, $\frac{du}{dy} + \frac{d\beta}{dx}$ we therefore have, setting the former $= X$ and the latter $= Y$, not merely, generally speaking, the equation $\frac{dX}{dx} + \frac{dY}{dy} = o$, but there is also extended through the whole boundary of any part of $T$
\[
\int \Bigl(X\frac{dx}{dp} + Y\frac{dy}{dp}\Bigr)\,ds = o, \quad \text{in so far as this expression at all has a determinate}
\]
value.

If we therefore decompose (according to Art.~9, V) the surface $T$, if it possesses a multiple connectivity, by crosscuts into a simply connected $T^*$, the integral
\[
\int_{O_0}^O \Bigl(\frac{du}{dp} - \frac{d\beta}{ds}\Bigr)\,ds \quad \text{for every line in the interior of } T^* \text{ going from } O_0 \text{ to } O \text{ has the same}
\]
value, and forms, $O_0$ being thought as fixed, a function of $x$, $y$ which suffers in $T^*$ everywhere a continuous change and along a crosscut an equal change on both sides. This function $\nu$ added to $\beta$ furnishes us a function $v = \beta + \nu$ whose differential quotient is $\frac{dv}{dx} = -\frac{du}{dy}$ and $\frac{dv}{dy} = \frac{du}{dx}$.

We therefore have the following

\emph{Theorem.} If in a connected surface $T$ decomposed by crosscuts into a simply connected $T^*$ there is given a complex function $\alpha + \beta i$ of $x$, $y$ for which
\[
\int \biggl[\Bigl(\frac{d\alpha}{dx} - \frac{d\beta}{dy}\Bigr)^2 + \Bigl(\frac{d\alpha}{dx} + \frac{d\beta}{dy}\Bigr)^2\biggr] dT \quad \text{extended through the whole surface has a finite}
\]\ednote{Thus in the print; $(d\alpha/dy + d\beta/dx)^2$ is presumably meant.}
value, it can always and only in one way be converted into a function of $z$ by addition of a function $\mu + \nu i$ of $x$, $y$ satisfying the following conditions:
1) $\mu$ is $= o$ at the boundary or at least different from it only at isolated points, $\nu$ arbitrarily given at a single point,


2) the changes of $\mu$ in $T$, of $\nu$ in $T^*$ are discontinuous only at isolated points and only in such a way that $\displaystyle\int \biggl[\Bigl(\frac{d\mu}{dx}\Bigr)^2 + \Bigl(\frac{d\mu}{dy}\Bigr)^2\biggr]\,dT$ and $\displaystyle\int \biggl[\Bigl(\frac{d\nu}{dx}\Bigr)^2 + \Bigl(\frac{d\nu}{dy}\Bigr)^2\biggr]\,dT$ extended through the whole surface remain finite, and the latter equal on both sides along the crosscuts.

The sufficiency of the conditions for the determination of $\mu + \nu i$ follows from the fact that $\mu$, by which $\nu$ is determined up to an additive constant, always at the same time yields a minimum of the integral $\Omega$, since, setting $u = \alpha + \mu$, evidently for each $\lambda$ $N = o$; a property that according to Art.~16 can belong to only one function.

\origpage{26}
\section*{19.}

The principles underlying the theorem at the conclusion of the preceding Art.\ open the way to investigate \emph{determinate} functions of a variable complex quantity (independently of an expression for the same).

For orientation in this field an estimate of the extent of the conditions required for the determination of such a function within a given domain of quantities will serve.

If we keep first to a determinate case, then, if the surface spread out over $A$ by which this domain of quantities is represented is a simply connected one, the function $w = u + vi$ of $z$ can be determined according to the following conditions:
1) for $u$ there is given at all boundary points a value that varies for an infinitely small change of position by an infinitely small quantity of the same order, but is otherwise arbitrary${}^{*)}$;


2) the value of $v$ is arbitrarily given at any point;


3) the function shall be finite and continuous at all points.
By these conditions, however, it is completely determined.

In fact, this follows from the theorem of the preceding Art., if one determines $\alpha + \beta i$, as will always be possible, such that $\alpha$ is equal at the boundary to the given value and throughout the entire surface the change of $\alpha + \beta i$ for every infinitely small change of position is infinitely small of the same order.

There can therefore, generally speaking, be given for $u$ at the boundary a completely arbitrary function of $s$, and thereby $v$ is co-determined everywhere; conversely, however, $v$ can also be assumed arbitrarily at each boundary point, from which the value of $u$ then follows. The scope for the choice of the values of $w$ at the boundary comprises therefore a manifold of one dimension for each boundary point, and their complete determination requires for each boundary point one equation, wherein it will nevertheless not be essential that each of these equations relates solely to the value of one term at one boundary point. This determination will also be able to happen in such a way that for each boundary point there is given one equation containing both terms and continuously changing its form with the position of this point, or for several parts of the boundary simultaneously in such a way that to each point of one of these parts there are associated $n - 1$ determinate points, one from each of the remaining parts, and for each $n$ such points jointly there are given $n$ equations continuously variable with their position. These conditions, whose totality forms a continuous manifold and which are expressed by equations between arbitrary functions, will nevertheless, in order to be admissible and sufficient for the determination of a function everywhere continuous in the interior of the domain of quantities, generally speaking, still require a restriction or supplementation

\textbf{*)} In themselves the changes of this value are subjected only to the restriction of not being discontinuous along a part of the boundary; a further restriction is made only in order to avoid unnecessary prolixity here.

\origpage{27}
by individual equations of condition --- equations for arbitrary constants ---, since the accuracy of our estimate evidently does not extend to these.

For the case where the domain of variability of the quantity $z$ is represented by a multiply connected surface, these considerations suffer no essential modification, in that the application of the theorem in Art.~18 yields a function constituted just as before except for the changes upon crossing the crosscuts --- changes that can be made $= o$ if the boundary conditions contain a number of available constants equal to the number of crosscuts.

The case where continuity in the interior along a line is dispensed with subordinates itself to the preceding if one regards this line as a cut of the surface.

If finally a violation of continuity, thus according to Art.~12 a becoming infinite of the function, is admitted at an isolated point, then, while retaining the other presuppositions made in our initial case, a function of $z$ after whose subtraction the function to be determined is to become continuous can be given arbitrarily for this point; by this, however, it is completely determined. For if one assumes the quantity $\alpha + \beta i$ in an arbitrarily small circle described about the point of discontinuity equal to this given function, but otherwise according to the former prescriptions, the integral
\[
\int \biggl[\Bigl(\frac{d\alpha}{dx} - \frac{d\beta}{dy}\Bigr)^2 + \Bigl(\frac{d\alpha}{dy} + \frac{d\beta}{dx}\Bigr)^2\biggr] dT \quad \text{extended over this circle } = o,
\]
and extended over the remaining part equal to a finite quantity, and one can thus apply the theorem of the preceding Art., whereby one obtains a function with the required properties. From this one can conclude with the aid of the theorem in Art.~13 that in general, when at an isolated point of discontinuity the function may become infinitely large of the order $n$, a number of $2n$ constants becomes available.

Geometrically represented, a function $w$ of a complex quantity $z$ variable within a given domain of quantities of two dimensions furnishes (according to Art.~15) from a given surface $T$ covering $A$ an image $S$ covering $B$ similar to it in the smallest parts, isolated points excepted. The conditions that have just been found sufficient and necessary for the determination of the function relate to its value either at boundary points or at points of discontinuity; they therefore all appear (Art.~15) as conditions for the position of the boundary of $S$, and indeed they give for each boundary point one equation of condition. If each relates only to one boundary point, they are represented by a family of curves, of which one forms the geometric locus for each boundary point. If two boundary points advancing continuously with each other are subjected jointly to two equations of condition, there arises thereby between two boundary parts such a dependence that, when the position of the one is arbitrarily assumed, the position of the other follows from it. In a similar manner a geometric meaning results for other forms of the equations of condition, which we do not wish to pursue further, however.

\section*{20.}

The introduction of complex quantities into mathematics has its origin and nearest

\origpage{28}
purpose in the theory of simple${}^{*)}$ laws of dependence between variable quantities expressed by operations on quantities. For if one applies these laws of dependence in an extended scope, by giving complex values to the variable quantities to which they relate, there emerges a harmony and regularity that otherwise remains concealed. The cases in which this has occurred comprise indeed as yet only a small domain --- they can almost all be traced back to those laws of dependence between two variable quantities where the one is either an algebraic${}^{**)}$ function of the other or such a function whose differential quotient is an algebraic function --- but almost every step that has been taken here has not merely given to the results obtained without the aid of complex quantities a simpler, more rounded shape, but has also broken the path to new discoveries, to which the history of the investigations on algebraic functions, circular or exponential functions, elliptic and Abelian functions furnishes the evidence.

It shall be briefly indicated what has been gained by our investigation for the theory of such functions.

The previous methods of treating these functions always laid at the foundation as definition an expression of the function, whereby its value was given for each value of its argument; by our investigation it has been shown that, in consequence of the general character of a function of a variable complex quantity, in a definition of this kind a part of the determining elements is a consequence of the remainder, and indeed the extent of the determining elements has been reduced to those necessary for determination. This essentially simplifies the treatment of the same. In order e.\ g.\ to prove the equality of two expressions of the same function, one formerly had to transform the one into the other, that is, show that both agreed for each value of the variable quantity; now the demonstration of their agreement in a far smaller extent suffices.

A theory of these functions on the foundations provided here would establish the configuration of the function (that is, its value for each value of its argument) independently of a mode of determining it by operations on quantities, by adding to the general concept of a function of a variable complex quantity only the characteristics necessary for the determination of the function, and only then proceed to the various expressions of which the function is capable. The common character of a genus of functions that are expressed in a similar manner by operations on quantities then presents itself in the form of the boundary and discontinuity conditions imposed on them. If e.\ g.\ the domain of variability of the quantity $z$ is extended simply or multiply over the entire infinite plane $A$, and within the same a discontinuity is permitted to the function only at isolated points, and indeed only a becoming infinite whose order is finite (where for an infinite $z$ this quantity itself, but for each finite value $z'$ of the same

\textbf{*)} We consider here as elementary operations addition and subtraction, multiplication and division, integration and differentiation, and a law of dependence as all the simpler, the fewer elementary operations condition the dependence. In fact, all functions hitherto used in analysis can be defined by a finite number of these operations.

\textbf{**)} That is, where an algebraic equation subsists between the two.

\origpage{29}
$\frac{1}{z - z'}$ counts as an infinitely large of the first order), then the function is necessarily algebraic, and conversely every algebraic function fulfills this condition.

The execution of this theory, which, as remarked, is destined to illuminate simple laws of dependence conditioned by operations on quantities, we omit for now, however, since we presently exclude the consideration of the expression of a function.

For the same reason we also do not occupy ourselves here with demonstrating the usefulness of our theorems as foundations of a \emph{general} theory of these laws of dependence, for which the proof is required that the concept of a function of a variable complex quantity laid at the foundation here completely coincides with that of a dependence expressible by operations on quantities${}^{*)}$.

\section*{21.}

An executed example of their application will nevertheless be of use for the illustration of our general theorems.

The application of them designated in the preceding article is, although the one primarily intended in their setting up, yet only a special one. For when the dependence is conditioned by a finite number of the operations on quantities considered there as elementary operations, the function contains only a finite number of parameters, which has the result for the form of a system of mutually independent boundary and discontinuity conditions sufficient for its determination that among them conditions to be arbitrarily determined at each point along a line cannot occur at all. For our present purpose it therefore seemed more suitable not to choose an example taken from there, but rather one where the function of the complex variable depends on an arbitrary function.

For visualization and more convenient expression we give to it the geometric clothing used at the conclusion of Art.~18. It then appears as an investigation into the possibility of furnishing from a given surface a connected image similar in the smallest parts whose shape is given, where thus, expressed in the above form, for each boundary point of the image a locus curve, and indeed for all the same one, and besides (Art.~5) the sense of the boundary and the winding points of the same are given. We restrict ourselves to the solution of this problem in the case where to each point of the one surface only one point of the other is to correspond and the surfaces are simply connected, for which case it is contained in the following theorem.

Two given simply connected plane surfaces can always be related to one another in such a way that to each point of the one there corresponds a point of the other advancing continuously with it and their corresponding smallest parts are similar; and indeed, for one interior

\textbf{*)} By this is understood every dependence expressible by a finite or infinite number of the four simplest operations of calculation, addition and subtraction, multiplication and division. The expression operations on quantities is intended to indicate (in contrast to operations on numbers) such operations of calculation in which the commensurability of the quantities does not come into consideration.

\origpage{30}
point and for one boundary point the corresponding ones can be arbitrarily given; by this, however, the relation is determined for all points.

If two surfaces $T$ and $R$ are related to a third $S$ in such a way that similarity obtains between the corresponding smallest parts, there results from this a relation between the surfaces $T$ and $R$ of which the same evidently holds. The problem of relating two arbitrary surfaces to one another in such a way that similarity obtains in the smallest parts is thereby reduced to that of mapping every arbitrary surface by a single determinate one similarly in the smallest parts. According to this, if we describe in the plane $B$ about the point where $w = o$ with radius $1$ a circle $K$, we need only prove in order to demonstrate our theorem: An arbitrary simply connected surface $T$ covering $A$ can always be mapped by the circle $K$ connectedly and similarly in the smallest parts, and indeed only in one way such that to the center there corresponds an arbitrarily given interior point $O_0$ and to an arbitrarily given point of the circumference an arbitrarily given boundary point $O'$ of the surface $T$.

We denote the determinate meanings of $z$, $Q$ for the points $O_0$, $O'$ by corresponding indices and describe in $T$ about $O_0$ as center an arbitrary circle $\varTheta$ that does not extend to the boundary of $T$ and contains no winding point. If we introduce polar coordinates by putting $z - z_0 = r e^{\varphi i}$, the function becomes $\log(z - z_0) = \log r + \varphi i$. The real value therefore varies continuously in the entire circle with the exception of the point $O_0$ where it becomes infinite. The imaginary, however, if everywhere the smallest positive among the possible values of $\varphi$ is chosen, receives along the radius where $z - z_0$ assumes real positive values the value $o$ on the one side, on the other the value $2\pi$, but then varies continuously at all remaining points. Evidently this radius can be replaced by a completely arbitrary line $l$ drawn from the center to the circumference, so that the function $\log(z - z_0)$ suffers a sudden decrease by $2\pi i$ upon the passage of the point $O$ from the negative (that is, where according to Art.~8 $p$ becomes negative) to the positive side of this line, but otherwise varies continuously with its position in the entire circle $\varTheta$. Let us now assume the complex function $\alpha + \beta i$ of $x$, $y$ in the circle $\varTheta = \log(z - z_0)$, but outside it, by extending $l$ arbitrarily up to the boundary, such that it
1) becomes at the circumference of $\varTheta = \log(z - z_0)$, and at the boundary of $T$ purely imaginary,


2) changes upon the passage from the negative to the positive side of the line $l$ by $-2\pi i$, but otherwise for each infinitely small change of position by an infinitely small quantity of the same order,
which will always be possible: then the $\displaystyle\int \biggl[\Bigl(\frac{d\alpha}{dx} - \frac{d\beta}{dy}\Bigr)^2 + \Bigl(\frac{d\alpha}{dy} + \frac{d\beta}{dx}\Bigr)^2\biggr]\,dT$ extended over $\varTheta$ receives the value zero, extended over the entire remaining part a finite value, and $\alpha + \beta i$ can therefore by addition of a continuous function of $x$, $y$, determined up to a purely imaginary constant

\origpage{31}
remainder, which is purely imaginary at the boundary, be converted into a function $t = m + ni$ of $z$. The real part $m$ of this function becomes at the boundary $= o$, at the point $O_0 = -\infty$, and varies continuously in the entire remainder of $T$. For each value $a$ of $m$ lying between $o$ and $-\infty$, $T$ therefore falls by a line where $m = a$ into parts where $m < a$ and which contain $O_0$ in the interior, on the one hand, and on the other hand into parts where $m > a$ and whose boundary is formed partly by the boundary of $T$, partly by lines where $m = a$. The order of connectivity of the surface $T$ is by this decomposition either not changed or lowered; the surface therefore falls, since this order is $= -1$, either into two pieces of order of connectivity $o$ and $-1$, or into more than two pieces. The latter, however, is impossible, because then in at least one of these pieces $m$ would have to be everywhere finite and continuous and constant in all parts of the boundary, and consequently would have to have either a constant value in a surface part, or somewhere --- at a point or along a line --- a maximum or minimum value, contrary to Art.~11, III. The points where $m$ is constant form therefore closed everywhere simple lines which bound a piece enclosing the point $O_0$, and indeed $m$ necessarily decreases towards the interior, from which it follows that during a positive circuit (where according to Art.~8 $s$ increases) $n$, as far as it is continuous, constantly increases, and thus, since it suffers only upon the passage from the negative to the positive side of the line $l$ a sudden change by $-2\pi^{*)}$, becomes equal to each value between $o$ and $2\pi$ once, apart from a multiple of $2\pi$. If we now put $e^t = w$, $e^m$ and $n$ become polar coordinates of the point $Q$ with respect to the center of the circle $K$. The totality of the points $Q$ then evidently forms a surface $S$ everywhere simply spread out over $K$; the point $Q_0$ of the same falls on the center of the circle; but the point $Q'$ can by means of the constant still available in $n$ be shifted to an arbitrarily given point of the circumference, q.\ e.\ d.\ednote{Thus in the print; presumably w.\ z.\ b.\ w. [q.\ e.\ d.] is meant.}

In the case where the point $O_0$ is a winding point of the $n-1$th order, one arrives at the goal, if only $\log(z - z_0)$ is replaced by $\frac{1}{n}\log(z - z_0)$, by entirely similar inferences, whose further execution one will easily supply, however, from Art.~14.

\section*{22.}

The complete execution of the investigation of the preceding article for the more general case where to one point of the one surface several points of the other are to correspond, and a simple connectivity is not presupposed for the same, we omit here, first of all

\textbf{*)} Since the line $l$ leads from a point situated in the interior of the piece up to an exterior one, it must, if it intersects its boundary multiple times, go once more from interior to exterior than from exterior to interior, and the sum of the sudden changes of $n$ during a positive circuit is therefore always $= 2 - \pi$\ednote{Thus in the print; presumably $-2\pi$ is meant.}.

\origpage{32}
since, conceived from a geometrical point of view, our entire investigation could have been conducted in a more general form. The restriction to plane surfaces, simple except for isolated points, is namely not essential for the same; rather, the problem of mapping an arbitrarily given surface similarly in the smallest parts onto another arbitrarily given surface admits of an entirely similar treatment. We content ourselves on this with referring to two \emph{Gauss}ian memoirs, that cited in Art.~3 and the disquis.\ gen.\ circa superf.\ art.~13.

\end{document}
